\chapter*{おわりに}\label{afterword}\addcontentsline{toc}{chapter}{おわりに}\markboth{おわりに}{おわりに}

ここまでで構文解析の世界を概観してみましたがいかがでしたか？構文解析、特に非自然言語の構文解析というのは地味なもので、パーサージェネレータの発展などもあり、20世紀末には「構文解析はもう終わった問題だ」という人もいました。その一方で2000年代以降になってもPEGの提案があり、Pythonの構文解析器に採用されるまでになりましたし、\texttt{LL(*)}や\texttt{ALL(*)}のような革新的なアルゴリズムが生み出されています。それも、主流であった上向き型の構文解析でなく下向き型の構文解析で、です。

とはいえ地味なものは地味であり、プログラミング言語処理系を構成するコンポーネントという観点からは「脇役」という印象は否めません。ただ、私たちはプログラミング言語を書いているときは、コンパイラの内部表現や抽象構文木と対話しているわけではありません。プログラマが直接対話する相手はプログラミング言語の具象構文であり、具象構文はプログラミング言語の「UI」を担当する部分といえるでしょう。通常のアプリケーションでUIが軽視されるべきでないのと同様にやはり具象構文も軽視されるべきでないと私は思いますし、よりよい具象構文の設計には構文解析の知識が助けになると信じています。

ここまで本書では、LL法やLR法といった古典的なアルゴリズムから、PEGとPackrat Parsing、構文解析器生成系、そして現実のプログラミング言語が抱える泥臭い問題までを見てきました。しかし、構文解析の世界は広く、本書で十分に扱えなかったテーマも数多く残っています。最後に、その一端を紹介しておきましょう。この本を読み終えた皆さんが次へ進むための、ささやかな地図になれば幸いです。

\begin{itemize}
\item
  \textgt{曖昧な文法とgeneralized parsing：}

  本書で扱った手法は、いずれも「1つの解析結果を効率よく求める」ことを前提としていました。これに対して、曖昧な文脈自由文法をそのまま受け取り、あり得る解析結果をすべて求める系譜があります。Earley法、GLR法、GLL法などがその代表で、generalized parsingと総称されます。第4章と第5章で名前だけ登場したGLR法も、元々は自然言語を解析するために生まれたものでした。プログラミング言語の世界では曖昧さは避けるべきものとして扱われますが、自然言語ではむしろ曖昧さこそが本質であり、それを正面から引き受けるための道具立てがこの系譜だ、と見ることもできます。プログラミング言語の枠を越えて構文解析を眺めてみたい方には、この系譜をおすすめします。参考文献の「Parsing Techniques」は、この世界まで含めて扱っている数少ない本です。

\item
  \textgt{インクリメンタル構文解析：}

  IDEやエディタは、キー入力のたびにファイル全体を解析し直しているわけにはいきません。編集された箇所だけを再解析して構文木を更新する技術がインクリメンタル構文解析です。本書では序章のコラムで触れるに留めましたが、構文解析器生成系tree-sitterのようにこの技術を核としたツールも普及しており、LSPが当たり前になった現在、実用上の重要性はますます高まっています。差分だけを解析し直すには、前回の解析結果をどう安全に再利用するかという、本書では扱わなかった種類の工夫が必要になります。

\item
  \textgt{エラーリカバリ：}

  第6章で駆け足で紹介しましたが、本来これだけで一冊の本が書ける主題です。正しい入力だけを相手にする構文解析器は比較的きれいに書けますが、IDEなどの開発環境が相手にするのは、ほとんどの場合「書きかけの壊れたプログラム」だからです。壊れた入力から、どれだけ意味のある構文木を回収できるかが、実用の構文解析器の品質では効いてきます。

\item
  \textgt{形式言語理論：}

  第3章で概観した正規言語や文脈自由言語といった言語クラスには、対応する計算モデル（機械）があることが知られています。正規言語は有限オートマトンで、文脈自由言語は有限オートマトンにスタックを追加した\textgt{プッシュダウンオートマトン}で認識できます。有限オートマトンで括弧の対応が扱えなかったのは、まさにこのスタックの有無の差なのです。さらに外側には、線形拘束オートマトンやチューリングマシンといった、より強力な計算モデルの世界が続いています。「なぜ\(a^n b^n\)は文脈自由言語なのに\(a^n b^n c^n\)はそうではないのか」といった問いに正面から答えたくなったら、形式言語理論の教科書に進むときです。

\end{itemize}

本書で学んだ文脈自由文法や各種アルゴリズムの理解が、これらのテーマに挑むときの土台になります。

以下では、こうしたテーマをさらに学ぶための文献を紹介します。

\section*{古典的名著・専門書}\label{a.1}\addcontentsline{toc}{section}{古典的名著・専門書}\markboth{おわりに}{古典的名著・専門書}

\begin{itemize}
\item
  A.V.エイホ、M.S.ラム、R.セシィ、J.D.ウルマン，『コンパイラ -
  原理・技法・ツール』（第2版）、サイエンス社、2009年（通称ドラゴンブック）

  \begin{itemize}
  \item
    コンパイラ構築に関する標準的な教科書です。字句解析、構文解析（LL、LR）、意味解析、コード生成など、コンパイラの全般的なトピックを網羅。理論的背景もしっかり解説されています。中級者以上向けです。
  \end{itemize}
\item
  J.ホップクロフト、R.モトワニ、J.ウルマン、『オートマトン 言語理論 計算論 Ⅰ』(第2版)
  、サイエンス社、2003年

  \begin{itemize}
  \item
    オートマトンと形式言語の理論に関する大学生レベル以上向け教科書。正規言語、文脈自由言語など、計算理論の基礎をしっかり学べます。数学的な厳密さを求める方向け。
  \end{itemize}
\item
  J.ホップクロフト、R.モトワニ、J.ウルマン、『オートマトン 言語理論 計算論 Ⅱ』(第2版)
  、サイエンス社、2003年

  \begin{itemize}
  \item
    同上。Ⅱではチューリングマシン、決定不能性、計算複雑性などを取り扱っています。Ⅱは計算理論のより深い部分に踏み込んでいます。
  \end{itemize}
\item
  Dick Grune, Ceriel J.H. Jacobs, \emph{Parsing Techniques: A Practical
  Guide} (2nd Edition), Springer, 2008年

  \begin{itemize}
  \item
    書名どおり、さまざまな構文解析技術に特化した書籍です。LL、LRだけでなく、アーリー法、GLR、CYK法など、高度なアルゴリズムや曖昧性のある文法の扱いについても詳しく解説されています。構文解析を専門的に極めたい方向け。
  \end{itemize}
\end{itemize}

\section*{特定の技術に関する論文・資料}\label{a.2}\addcontentsline{toc}{section}{特定の技術に関する論文・資料}\markboth{おわりに}{特定の技術に関する論文・資料}

\begin{itemize}
\item
  Bryan Ford, ``Parsing Expression Grammars: A Recognition-Based
  Syntactic Foundation'', 2004

  \begin{itemize}
  \item
    PEGを提案したオリジナルの論文です。PEGの形式的な定義、操作的意味論などについて説明があります。理論的な背景を深く理解したい方向け。
  \end{itemize}
\item
  Tim A. Wagner, Susan L. Graham, ``Efficient and Flexible Incremental
  Parsing'', ACM Transactions on Programming Languages and Systems,
  20(5), 1998

  \begin{itemize}
  \item
    インクリメンタル構文解析の代表的な論文です。編集された箇所だけを安全に再解析するための理論と実装を詳述しており、tree-sitterをはじめとする現代のツールの基礎となっています。
  \end{itemize}
\item
  Lukas Diekmann, Laurence Tratt, ``Don't Panic! Better, Fewer, Syntax
  Errors for LR Parsers'', ECOOP 2020

  \begin{itemize}
  \item
    LR構文解析のエラーリカバリを実用的な品質へ引き上げた比較的新しい論文です。エラーリカバリ研究の現在地を知る入口として適しています。オープンアクセスで公開されています。
  \end{itemize}
\item
  Terence Parr, \emph{The Definitive ANTLR 4 Reference}, Pragmatic
  Bookshelf (2nd Edition), 2013年

  \begin{itemize}
  \item
    ANTLR v4の作者自身による解説書です。ANTLRの文法定義、使い方、\texttt{ALL(*)}アルゴリズムの概要、実践的なパーサー構築のテクニックが豊富。ANTLRを使いこなしたいなら必読。
  \end{itemize}
\item
  ANTLR公式サイト、\url{https://www.antlr.org/}

  \begin{itemize}
  \item
    ANTLRのドキュメント、チュートリアル、文法リポジトリなど。最新情報やコミュニティのサポートも得られます。
  \end{itemize}
\item
  JavaCC公式サイト、\url{https://javacc.github.io/javacc/}

  \begin{itemize}
  \item
    JavaCCのドキュメント、チュートリアル、FAQなど。
  \end{itemize}
\item
  GNU Bisonマニュアル、\url{https://www.gnu.org/software/bison/manual/}

  \begin{itemize}
  \item
    Bison (Yacc互換)の詳細なマニュアル。LALR(1)やGLRパーサーの生成方法、文法定義の書き方などが学べます。
  \end{itemize}
\item
  tree-sitter公式サイト、\url{https://tree-sitter.github.io/tree-sitter/}

  \begin{itemize}
  \item
    インクリメンタル構文解析とエラーリカバリを核とした、エディタ向けの構文解析器生成系です。多くのエディタやコード解析ツールで採用されています。
  \end{itemize}
\end{itemize}

これらの資料を通じて構文解析の世界への探求をさらに深めていただければ幸いです。
